\section{Introduction}\label{sec:intro}

A professional basketball career is a sequence, not a number: two players with the same career-average productivity may have reached it through an early peak and a long decline, or through a slow rise that culminates late. Most quantitative work on National Basketball Association (NBA) players nevertheless summarises a career by aggregate statistics, or models the population-average shape of the age--performance relationship \citep{vaci2019,kalen2021,xiao2025}. Whether individual careers fall into a small number of qualitatively distinct trajectory \emph{types}, or instead vary continuously around a common curve, is an open empirical question with consequences for how career data should be modelled and how career-based comparisons between players should be made.

Unsupervised learning offers an apparently direct route to that question: represent every career as a time series, learn a compact representation, and cluster. Recent pipelines that combine a learned embedding, nonlinear dimensionality reduction with UMAP \citep{mcinnes2018} and density-based clustering with HDBSCAN \citep{campello2013,campello2015} have become a default recipe for structure discovery in many fields \citep{angelov2020,grootendorst2022,fragkedaki2024}. The recipe requires neither a fixed number of clusters nor a parametric trajectory model, but it is also fragile: the partition can depend on random seeds, on the subsample and on interacting hyper-parameters, and internal validity indices computed in the clustering space can reward artefacts of the embedding rather than structure in the data \citep{chari2023,vonluxburg2010}. In an earlier version of this work we applied such a pipeline to Player Efficiency Rating (PER) sequences and reported nineteen trajectory clusters; editorial review asked whether that solution would survive a change of seed, latent dimensionality or sample restriction, whether it outperformed simpler methods, and whether the clusters meant anything beyond the indices used to select them. The present paper is a new study built around those questions, with new data construction, a new validation design and, as it turns out, a different substantive conclusion.

We treat the question of career-trajectory types as a problem of \emph{reproducible} structure discovery. Our contribution is methodological and empirical at the same time. Methodologically, we embed the representation-learning pipeline in a validation protocol that (i) separates model selection from evaluation through a held-out set of players, (ii) quantifies the reproducibility of every partition across autoencoder seeds, UMAP seeds, subsamples and hyper-parameters, (iii) replaces single-run solutions by resampling-based consensus clustering \citep{monti2003,fred2005} with a prespecified rule for the number of clusters, (iv) compares the learned representation with summary-statistic, principal-component and dynamic-time-warping representations under identical clustering and evaluation rules, (v) validates the resulting types against variables that were not used in clustering, and (vi) tests whether reproducible partitions reflect discrete structure at all, against structureless reference data, offering continuous career coordinates as the alternative summary when they do not. Empirically, we apply the protocol to \nPlayers{} players who entered the NBA between the 1979--80 and 2010--11 seasons and whose careers are therefore complete or nearly complete, using PER as computed by Basketball-Reference and, as a robustness check, Box Plus-Minus (BPM).

The study is organised around five research questions with success criteria fixed before the analyses were run; the one criterion that had to be changed during the analysis, the rule for the number of consensus clusters, is documented in Section~\ref{sec:consensus}.
\begin{description}
\item[RQ1] How accurately does a Transformer autoencoder reconstruct variable-length PER careers on held-out players, compared with a global-mean, a player-mean and a player-specific linear-trend baseline, and how does accuracy depend on latent dimensionality? \emph{Criterion:} the autoencoder is retained only if its held-out root-mean-square error (RMSE) is below that of the linear-trend baseline.
\item[RQ2] Is the partition produced by a single UMAP+HDBSCAN run reproducible? \emph{Criterion:} a design is called reproducible if the mean adjusted Rand index (ARI) between runs that differ only in random seed is at least $0.80$.
\item[RQ3] Which number of trajectory types is supported by consensus clustering, and what are their profiles? \emph{Criterion:} the number of clusters is chosen by a fixed rule on the consensus co-association matrix (Section~\ref{sec:consensus}), and a partition is interpreted as a typology only if it also passes tests against structureless reference data.
\item[RQ4] Does the learned representation improve on summary statistics, principal components and dynamic time warping in reproducibility, trajectory-shape homogeneity and external validity when all are clustered with the same consensus protocol?
\item[RQ5] Are the types associated with variables not used in clustering (position, draft status, All-Star and All-NBA selection, Win Shares, BPM), and are they robust to sample restrictions, to latent dimensionality and to the choice of performance metric?
\end{description}

Code, the processed dataset and every intermediate result are released with the paper.
